\documentclass[12pt]{article}
\usepackage[utf8]{inputenc}
\usepackage[T1]{fontenc}
\usepackage{hyperref}
\title{Coyutla Totonac Language Resources for K-12 Curriculum Development}
\author{Compiled by Claude Sonnet 5 (claude-sonnet-5) for Andr\'as Kornai}
\date{}
\begin{document}
\maketitle

\section{Introduction}
Coyutla Totonac (ISO 639-3: toc; Glottocode: coyu1235) is a Totonacan language of Mexico, spoken in Puebla and Veracruz. This search found very little dedicated material, and the list below is short.

\section{Primary Educational Resources}
The search found no primers, textbooks or school materials for Coyutla Totonac.

\section{Linguistic and Reference Materials}
Glosbe has user-contributed dictionary pages pairing Coyutla Totonac with English and other languages. I did not open them, so their size and quality are unchecked: \url{https://glosbe.com/toc/en} and \url{https://glosbe.com/en/toc}.

\section{Religious and Community Resources}
The search found nothing.

\section{Digital and Audio Resources}
Wikitongues has a page with one video sample of the language: \url{https://wikitongues.org/languages/toc/}.

\section{Curriculum Development Context}
The search found no evidence of school use, no grammar and no published dictionary. A curriculum would have to rely on new documentation. Academic and Mexican institutional sources (for example INALI) still need a follow-up search.
\end{document}